\section{Anomaly Scoring Framework}
\label{sec:scoring}

We evaluate four families of anomaly scoring functions, each capturing a different aspect of reconstruction mismatch.
Let $\mathbf{y} = (y_1, \ldots, y_T)$ denote the claimed G-code token sequence aligned to a sensor window, and let $p(y_t \mid y_{<t}, \tilde{\mathbf{z}})$ denote the decoder's conditional probability for token $y_t$ given preceding tokens and the operation-conditioned encoder memory $\tilde{\mathbf{z}}$.

\subsection{Negative Log-Likelihood (NLL) Scores}
\label{sec:nll_score}

NLL-based scores measure decoder ``surprise'' when presented with the claimed G-code sequence.
Under normal operation, the decoder assigns high probability to correct tokens (low NLL); under attack, predictions diverge from claimed tokens (elevated NLL).

\paragraph{Mean NLL ($S_{\text{mean}}$).}
The mean NLL averages per-token negative log-probability across all $T$ tokens:
\begin{equation}
    S_{\text{mean}}(\mathbf{y}) = -\frac{1}{T} \sum_{t=1}^{T} \log p(y_t \mid y_{<t}, \tilde{\mathbf{z}})
    \label{eq:s_mean}
\end{equation}
This is equivalent to the log-perplexity: $\text{PPL}(\mathbf{y}) = \exp(S_{\text{mean}})$.

\paragraph{Max-token NLL ($S_{\max}$).}
$S_{\max}$ selects the single most anomalous token:
\begin{equation}
    S_{\max}(\mathbf{y}) = \max_{t \in \{1, \ldots, T\}} \left[ -\log p(y_t \mid y_{<t}, \tilde{\mathbf{z}}) \right]
    \label{eq:s_max}
\end{equation}
This score is critical for the V8 multi-line decoder, where most tokens in an attacked window may be unmodified.
If an A2 perturbation modifies a single X value within a 63-token sequence, $S_{\text{mean}}$ dilutes the signal by averaging the anomalous token with 62 normal tokens, while $S_{\max}$ captures it directly.

\paragraph{Scale dependence.}
For V7 ($T = 3$ content tokens), $S_{\text{mean}}$ and $S_{\max}$ are highly correlated because a single anomalous token dominates both.
For V8 ($T \leq 63$), they diverge: $S_{\max}$ outperforms on localized attacks (e.g., single coordinate perturbation), while $S_{\text{mean}}$ outperforms on diffuse attacks (e.g., A4 cross-operation swap).

\subsection{Rank-Based Score}
\label{sec:rank_score}

The rank-based score measures how far the claimed token falls from the decoder's top prediction:
\begin{equation}
    S_{\text{rank}}(\mathbf{y}) = \frac{1}{T} \sum_{t=1}^{T} \frac{\text{rank}(y_t)}{|\mathcal{V}|}
    \label{eq:s_rank}
\end{equation}
where $\text{rank}(y_t)$ is the 1-indexed position of the claimed token in the decoder's probability-sorted vocabulary $\mathcal{V}$ at position $t$, and $|\mathcal{V}| = 712$.

Under normal operation, the correct token is typically ranked first or near the top ($S_{\text{rank}} \approx 0.001$).
Under attack, even when the perturbed value receives non-negligible probability (e.g., for small coordinate perturbations where the attacked value is numerically close to the original), its \emph{rank} drops because the decoder preferentially assigns higher probability to the original value.

\paragraph{Advantage over NLL.}
Rank-based scoring is \emph{invariant to absolute probability magnitudes}.
For low-confidence numeric tokens (Section~\ref{sec:results_two_populations}), the decoder assigns similar, modest probability to the original and perturbed tokens, so $\Delta\text{NLL}$ is small.
However, the decoder's \emph{ranking} may still favor the original value: if the original has rank 45 and the perturbed token rank 120, the rank score detects the shift despite negligible absolute probabilities.
Empirically, rank-based scoring outperforms NLL by up to 16 percentage points: on A2 with $\delta = 0.001$~mm, V7 Rank achieves AUROC = 0.859 versus V7 NLL = 0.697 (Section~\ref{sec:results_anomaly}).

\subsection{Confidence-Weighted Score}
\label{sec:confidence_score}

To exploit the two-population structure (Section~\ref{sec:results_two_populations}), we weight each token's NLL by the decoder's maximum predicted probability $w_t = \max_v p(v \mid y_{<t}, \tilde{\mathbf{z}})$:
\begin{equation}
    S_{\text{conf}}(\mathbf{y}) = \frac{\sum_{t=1}^{T} w_t \cdot \left[ -\log p(y_t \mid y_{<t}, \tilde{\mathbf{z}}) \right]}{\sum_{t=1}^{T} w_t}
    \label{eq:s_conf}
\end{equation}
Confident tokens (54.3\% of numeric positions on the axis-conditional confidence measure of Section~\ref{sec:results_two_populations}) produce strong signals; the lowest-confidence tokens contribute minimally.

\subsection{Multi-Head Disagreement Score}
\label{sec:disagreement_score}

We evaluate a disagreement score measuring inconsistency across the decoder's output heads:
\begin{equation}
    S_{\text{disagree}}(\mathbf{y}) = \frac{1}{T} \sum_{t=1}^{T} \mathbb{1}\left[ \hat{y}_t^{(\text{type})} \neq \text{expected\_type}(\hat{y}_t^{(\text{legacy})}) \right]
    \label{eq:s_disagree}
\end{equation}
where $\hat{y}_t^{(\text{type})}$ and $\hat{y}_t^{(\text{legacy})}$ are the argmax predictions of the type and legacy heads, and $\text{expected\_type}(\cdot)$ maps a vocabulary token to its grammatical type.

\paragraph{Negative result.}
Multi-head disagreement achieves AUROC $= 0.500$ across all attack types and perturbation magnitudes in both decoders---no better than chance.
The jointly trained heads share the same Transformer trunk and produce highly correlated outputs: under attack, all heads are confused in the same direction, producing consistent but incorrect predictions.
The baseline disagreement rate is only 0.008, providing insufficient variance for discrimination.
We report this because multi-head disagreement has been proposed as a general-purpose uncertainty signal; our results show it is ineffective when heads share a common backbone.

\subsection{Hybrid Multi-Scale Score}
\label{sec:hybrid_score}

Different scoring methods and decoder scales excel on different attack types (Section~\ref{sec:results_anomaly}), motivating a hybrid score:
\begin{equation}
    S_{\text{hybrid}} = \begin{cases}
        S_{\text{mean}}^{(\text{V7})} & \text{for structural attacks (A1 command swap, A3 parameter substitution)} \\
        S_{\text{rank}}^{(\text{V7})} & \text{for small coordinate perturbations (A2, $\delta \leq 0.1$~mm)} \\
        S_{\max}^{(\text{V8})} & \text{for large coordinate perturbations (A2, $\delta > 0.1$~mm)} \\
        S_{\text{rank}}^{(\text{V8})} & \text{for semantic/cross-operation attacks (A4)}
    \end{cases}
    \label{eq:s_hybrid}
\end{equation}

When the attack type is unknown, the hybrid score can be approximated by computing all component scores, z-score normalizing each using training-set statistics, and taking the maximum:
\begin{equation}
    S_{\text{hybrid-max}} = \max\left\{ \tilde{S}_{\text{mean}}^{(\text{V7})}, \; \tilde{S}_{\text{rank}}^{(\text{V7})}, \; \tilde{S}_{\max}^{(\text{V8})}, \; \tilde{S}_{\text{rank}}^{(\text{V8})} \right\}
    \label{eq:s_hybrid_max}
\end{equation}
where $\tilde{S}$ denotes z-score normalization.
This max-over-normalized-scores approach trades per-attack-type optimality for generality, at the cost of a slightly elevated false positive rate due to the heavier right tail of the maximum over multiple scores.
We did not separately evaluate hybrid-max; the grid-searched ensemble (Section~\ref{sec:results_anomaly_head}) provides a more principled weight optimization over the same component scores.
